\documentclass[12pt]{article}
\usepackage{amsmath, amsthm, amssymb, fullpage, amscd}
\usepackage[english]{babel}
\usepackage[utf8]{inputenc}
\usepackage{eucal}
\usepackage{thmtools}
\usepackage{float}
\usepackage{thm-restate}
\usepackage{tikz-cd, pgf, tikz}
\usepackage{geometry}\geometry{margin=1in}
\usepackage[colorlinks, bookmarks=true]{hyperref}
\hypersetup{
colorlinks, citecolor=blue, linkcolor=blue, urlcolor=blue}
\usepackage{cleveref, autonum}
\usepackage{titlesec}
\usepackage{xcolor}
\usepackage{mathrsfs}
\usepackage{mathtools}
\usepackage[nottoc]{tocbibind}
\usepackage[noadjust]{cite}
\usepackage{graphicx}
\usepackage{caption}
\usepackage{subcaption}
\usepackage{bm, bbm}
\usepackage{mathrsfs}
\usepackage{comment}
\usepackage{array}
\usepackage{latexsym}
\usepackage{stmaryrd}
\usepackage{commath,enumitem,chngcntr}
\usepackage{indentfirst, hyperref}
\usepackage[titletoc,toc,title]{appendix}
\usepackage{pst-node}
\usepackage{yhmath}
\usepackage{appendix}
\usepackage[inkscapelatex=false]{svg}
\usepackage{algorithm}
\usepackage[noend]{algpseudocode}

\usepackage{faktor}
% \mathtoolsset{showonlyrefs}

\usepackage{hyperref}

\setlist{itemsep=3pt plus 1pt minus 1pt}

\usetikzlibrary{arrows,calc,positioning}

%%%%%%%%%%%%%%%%%%%%%%%%%%%%%%%%%%%%%%%%%%

\renewcommand{\norm}[1]{\left\lVert#1\right\rVert}

\newcommand{\dmo}{\DeclareMathOperator}
\dmo{\Id}{Id}
\dmo{\pt}{pt}
\dmo{\ab}{ab}
\dmo{\GL}{GL}
\dmo{\SO}{SO}
\dmo{\SU}{SU}
\dmo{\SL}{SL}
\dmo{\PSL}{PSL}
\dmo{\im}{Im}
\dmo{\re}{Re}
\dmo{\Ad}{Ad}
\dmo{\ad}{ad}
\dmo{\tr}{tr}
\dmo{\Log}{Log}
\dmo{\Homeo}{Homeo}
\dmo{\Vol}{Vol}
\dmo{\supp}{supp}
\dmo{\Isom}{Isom}
\dmo{\Hom}{Hom}
\dmo{\Ext}{Ext}
\dmo{\Tor}{Tor}
\dmo{\ind}{ind}
\dmo{\PP}{\mathbf{P}}

%%%%%%%%%%%%%%%%%%%%%%%%%%%%%%%%%%%%%%%%%
\newcommand{\C}{\mathbb C}
\newcommand{\R}{\mathbb R}
\newcommand{\Q}{\mathbb Q}
\newcommand{\Z}{\mathbb Z}
\newcommand{\N}{\mathbb N}
\newcommand{\T}{\mathcal T}
\newcommand{\D}{\mathbb D}
\renewcommand{\H}{\mathbb H}

\newcommand{\ra}{\rightarrow}
\newcommand{\ep}{\epsilon}
\newcommand{\de}{\delta}
\newcommand{\al}{\alpha}
\newcommand{\be}{\beta}
\newcommand{\til}{\tilde}
\newcommand{\vn}{\varnothing}
\newcommand{\ga}{\gamma}
\newcommand{\Si}{\Sigma}
\newcommand{\De}{\Delta}
\newcommand{\si}{\sigma}
\newcommand{\pa}{\partial}
\newcommand{\ti}{\times}
\newcommand{\ld}{\ldots}
\newcommand{\cd}{\cdots}
\newcommand{\sbs}{\subset}
\newcommand{\xra}{\xrightarrow}
\newcommand{\hra}{\hookrightarrow}
\newcommand{\Arg}{\mathrm {Arg}}
\renewcommand{\tr}{\mathrm {tr}}
\newcommand{\lan}{\langle}
\newcommand{\ran}{\rangle}
\newcommand{\from}{\colon}
\newcommand{\parallelogram}{%
  \tikz[baseline=-0.5ex, scale=0.15]{
    \draw (0,-0.6) -- (2,-0.6) -- (3,1.4) -- (1,1.4) -- cycle;
  }%
}

\newtheorem{theorem}{Theorem}[section]
\newtheorem{corollary}[theorem]{Corollary}
\newtheorem{proposition}[theorem]{Proposition}
\newtheorem{lemma}[theorem]{Lemma}
\newtheorem{conjecture}[theorem]{Conjecture}
\newtheorem{question}[theorem]{Question}
\newtheorem{problem}[theorem]{Problem}

\theoremstyle{definition}
\newtheorem{definition}[theorem]{Definition}
\newtheorem{definitions}[theorem]{Definitions}
\newtheorem{example}[theorem]{Example}
\newtheorem{examples}[theorem]{Examples}
\newtheorem{notation}[theorem]{Notation}
\newtheorem{notations}[theorem]{Notations}

\theoremstyle{remark}
\newtheorem{remark}[theorem]{Remark}
\newtheorem{remarks}[theorem]{Remarks}
% \newtheorem{warning}[theorem]{Warning}

\makeatletter
\let\c@equation\c@theorem
\makeatother
\numberwithin{equation}{section}

\newcommand\numberthis{\refstepcounter{equation}\tag{\theequation}}

\usepackage{lipsum}
% ADD THE FOLLOWING COUPLE LINES INTO YOUR PREAMBLE
\let\OLDthebibliography\thebibliography
\renewcommand\thebibliography[1]{
	\OLDthebibliography{#1}
	\setlength{\parskip}{2pt}
	\setlength{\itemsep}{0pt plus 0.3ex}
}

\newcommand\mymark{\indent $\triangleright$\ \ }
\def\gl{\lambda}

\author{Yanwen Luo, Yuan Luo}

\title{Morphing Graphs on High Genus Surfaces}
\date{May 2026}

\begin{document}

\maketitle
\providecommand{\MorphFigurePath}{pictures}
\svgpath{{\MorphFigurePath/}}

\section{Experiments}
\label{sec:experiments}

We evaluate the proposed hyperbolic morphing procedure on four cut-open surface embeddings: a genus--2 triangulation, a genus--2 cellular graph with quadrilateral and pentagonal faces, a genus--3 triangulation, and a genus--3 cellular graph with quadrilateral and pentagonal faces. All embeddings are represented in the Poincaré disk. Side-pairing Möbius transformations encode the side identifications of the fundamental polygon. The genus--2 examples use a regular octagonal fundamental domain for the Bolza surface, while the genus--3 examples use a regular 14-gonal fundamental domain for the Klein quartic.

The combinatorial sizes of the four test cases are listed in Table~\ref{tab:morph-topology-counts}. We omit the construction details of the initial planar graphs here; for the purposes of the morphing experiment, the important distinction is whether the input is a triangulation or a cellular decomposition with higher-order faces.

\begin{table}[t]
\centering
\caption{Topology of the four morphing experiments.}
\label{tab:morph-topology-counts}
\begin{tabular}{lrrrrrr}
\hline
Example & Genus & Vertices & Edges & Triangles & Quads & Pentagons \\
\hline
Triangulation & 2 & 190 & 527 & 338 & 0 & 0 \\
Cells & 2 & 190 & 490 & 270 & 25 & 6 \\
Triangulation & 3 & 250 & 677 & 428 & 0 & 0 \\
Cells & 3 & 250 & 557 & 208 & 80 & 20 \\
\hline
\end{tabular}
\end{table}

\paragraph{Notation and terminology.}
The graph $G=(V,E)$ is combinatorial: it stores vertices, edges, and faces, but not the positions on the Poincaré disk. The position array $Z$ stores the positions of all vertices in the cut-open fundamental polygon. For vertices attached by Möbius transformations, we choose one copy as the representative vertex and call it the \emph{root}; the remaining attached copies are called \emph{slave} vertices. More explicitly, if $z_i^1$ is the root copy and $z_i^\ell$ is a slave copy, then the attaching map $M_\ell$ satisfies
\[
z_i^\ell = M_\ell(z_i^1), \qquad \ell=2,\ldots,m_i .
\]
Thus root and slave copies are both stored in $Z$, but only the root is an independent position variable. An interior root has no slave copies ($m_i=1$). A boundary root may have one or more slave copies attached by Möbius transformations. A fixed root is a polygon corner whose position is held fixed. A free root is any non-fixed root, including ordinary interior roots and non-corner boundary roots. Below we use only the implementation terms root, slave, free root, fixed root, and interior root.

\paragraph{Endpoint embeddings.}
In applications, the endpoint embeddings may be prescribed. In these experiments, we instead construct synthetic endpoints in order to compare the four graph types under the same procedure.

For each graph, we compute two endpoints by minimizing a hyperbolic Dirichlet energy. The source endpoint uses scalar edge weights $w_e \equiv 1$. The target endpoint uses a non-uniform scalar weight field in which the first half of the edge list has weight $10$, and the remaining half has weight $1$. Both endpoints are computed with the hyperbolic harmonic-map solver in the Poincaré disk, subject to the same fixed fundamental-domain corner vertices and the same side-pairing Möbius constraints.

\paragraph{Gradient computation.}
We write a vertex position as a complex number $z_i \in \D$. The Poincaré disk metric is conformal to the Euclidean metric with factor
\[
\lambda_z = \frac{2}{1-|z|^2}.
\]
For $x,y\in \D$, Möbius addition is
\[
x \oplus y = \frac{x+y}{1+\overline{x}y},
\]
and the hyperbolic distance is
\[
d_{\D}(x,y) = 2 \operatorname{arctanh}\left|(-x)\oplus y\right|.
\]
The logarithm and exponential maps used in the gradient descent are
\[
\log_x(y) = \frac{2}{\lambda_x} \frac{\operatorname{arctanh}|\delta|}{|\delta|} \delta, \qquad \delta=(-x)\oplus y,
\]
with $\log_x(x)=0$, and
\[
\exp_x(v) = x \oplus \left( \tanh\left(\frac{\lambda_x |v|}{2}\right) \frac{v}{|v|} \right),
\]
with $\exp_x(0)=x$.

For scalar endpoint weights, the energy is
\[
E(z)=\frac12\sum_{\{i,j\}\in E} w_{ij}\,d_{\D}(z_i,z_j)^2.
\]
The Riemannian gradient contribution at the two endpoints of an edge is
\[
\nabla_{z_i}E_{\{i,j\}}=-w_{ij}\log_{z_i}(z_j), \qquad \nabla_{z_j}E_{\{i,j\}}=-w_{ij}\log_{z_j}(z_i).
\]
For directed mean-value weights, we use directed coefficients $\omega_{ij}$ and $\omega_{ji}$ in the gradient. With $d_{ij}=d_{\D}(z_i,z_j)$, the edge contribution is
\begin{equation}
\label{eq:directed-edge-gradient}
\begin{aligned}
\nabla_{z_i}E_{\{i,j\}} &= -\omega_{ij}\frac{\sinh d_{ij}}{d_{ij}}\log_{z_i}(z_j), \\
\nabla_{z_j}E_{\{i,j\}} &= -\omega_{ji}\frac{\sinh d_{ij}}{d_{ij}}\log_{z_j}(z_i).
\end{aligned}
\end{equation}
We use Eq.~\eqref{eq:directed-edge-gradient} in the gradient-accumulation algorithm below, with $\sinh d_{ij}/d_{ij}$ set to $1$ when $d_{ij}=0$.

The factor $\sinh d_{ij}/d_{ij}$ is not part of the logarithm map and is not used for the scalar Dirichlet endpoint energy above. It appears when the directed mean-value solve is written using the cosh-distance edge energy $\omega_{ij}(\cosh d_{ij}-1)$. Since $d(\cosh x-1)/d x=\sinh x$ and
\[
\nabla_{z_i} d_{\D}(z_i,z_j) = -\frac{1}{d_{ij}}\log_{z_i}(z_j),
\]
the directed contribution becomes
\[
\nabla_{z_i}E_{\{i,j\}} = -\omega_{ij} \frac{\sinh d_{ij}}{d_{ij}} \log_{z_i}(z_j).
\]
Thus the factor comes from differentiating the cosh-distance energy. In contrast, the squared-distance Dirichlet energy gives simply $-w_{ij}\log_{z_i}(z_j)$. This distinction is useful for the mean-value weights because $(\sinh d_{ij}/d_{ij})\log_{z_i}(z_j)=\sinh d_{ij}u_{ij}$, where $u_{ij}$ is the unit tangent direction from $z_i$ to $z_j$.

\paragraph{Gradient update.}
The side-pairing constraint between a root and one of its slaves has the form
\[
z_s=M_{s\leftarrow r}(z_r),
\]
where $r$ is the root and $s$ is the slave copy. The update of each free root is performed by the exponential map,
\[
z_i^{(m+1)} = \exp_{z_i^{(m)}}\left(-\eta\,\nabla_{z_i}E\right),
\]
while fixed roots remain fixed. An interior root is just a free root with no slaves, so its accumulated gradient is computed directly from its incident edges and then used in the exponential-map update. A boundary root first collects gradient contributions from its root copy and from its slave copies; the slave contributions are pulled back to the root before the same root update is applied. Slave copies are not independently updated; after each free root update, the corresponding slave entries in $Z$ are overwritten by their side-pairing Möbius images. Algorithm~\ref{alg:root-slave-gradient-accumulation} gives the update used in the implementation.

\paragraph{Gradient transport across root and slave copies.}
Gradient contributions are computed on each copy in the cut-open polygon. Contributions from slave copies are transported to the root copy before applying the root update.

Each side-pairing map is an orientation-preserving disk isometry
\[
M(z)=\frac{az+b}{\overline{b}z+\overline{a}}, \qquad |a|^2-|b|^2=1.
\]
Its complex derivative is
\[
M'(z)=\frac{1}{(\overline{b}z+\overline{a})^2}.
\]
Let $z_s=M(z_r)$, and let $g_s\in T_{z_s}\D$ be the gradient contribution at the slave copy. The corresponding contribution at the root is the metric adjoint
\[
g_{r\leftarrow s} = (dM_{z_r})^{*}g_s.
\]
Using the conformal form of the Poincaré metric,
\[
\langle u,v\rangle_z=\lambda_z^2\operatorname{Re}(u\overline{v}),
\]
we obtain
\[
(dM_{z_r})^{*}g_s = \frac{\lambda_{z_s}^2}{\lambda_{z_r}^2} \overline{M'(z_r)}\,g_s.
\]
Because $M$ is a hyperbolic isometry, this simplifies to the inverse differential:
\[
(dM_{z_r})^{*}g_s = d(M^{-1})_{z_s}(g_s) = (M^{-1})'(z_s)\,g_s = \frac{g_s}{M'(z_r)}.
\]
Computationally, for $M(z)=(az+b)/(\overline{b}z+\overline{a})$, the inverse is
\[
M^{-1}(z)=\frac{\overline{a}z-b}{-\overline{b}z+a},
\]
and its derivative is
\[
(M^{-1})'(z)=\frac{1}{(-\overline{b}z+a)^2}.
\]
Thus the pulled-back gradient contribution used in the implementation is
\begin{equation}
\label{eq:pullback-gradient}
g_{r\leftarrow s} = (M^{-1})'(z_s)g_s = \frac{g_s}{(-\overline{b}z_s+a)^2}.
\end{equation}
If several slave copies attach to the same root, their pulled-back contributions are summed before the root update. Although $g_i^\ell$ is stored as a complex number, it represents a vector in $T_{z_i^\ell}\D$. Thus two equal complex numbers based at different copies are not the same tangent vector at the root until they are identified by the differential of the attaching Möbius map.

\begin{algorithm}[H]
\caption{Gradient accumulation and update at one root}
\label{alg:root-slave-gradient-accumulation}
\begin{algorithmic}[1]
\State \textbf{Input:} Root copy $z_i^1$, slave copies $z_i^2,\ldots,z_i^m$, attaching maps $M_\ell$ satisfying $z_i^\ell=M_\ell(z_i^1)$, directed weights $\omega$, and step size $\eta$
\State \textbf{Output:} Updated root and slave copies for vertex $i$
\For{$\ell=2,\ldots,m$}
    \State Enforce the slave constraint: $z_i^\ell\gets M_\ell(z_i^1)$
\EndFor
\State $G_i\gets 0$
\For{$\ell=1,\ldots,m$}
    \State Compute $g_i^\ell$ at copy $z_i^\ell$ from Eq.~\eqref{eq:directed-edge-gradient}
    \If{$\ell=1$}
        \State $G_i\gets G_i+g_i^\ell$ \Comment{root-copy contribution}
    \Else
        \State Pull back the slave contribution by Eq.~\eqref{eq:pullback-gradient}: $\widetilde g_i^\ell\gets d(M_\ell^{-1})_{z_i^\ell}(g_i^\ell)$
        \State $G_i\gets G_i+\widetilde g_i^\ell$
    \EndIf
\EndFor
\If{$z_i^1$ is a free root}
    \State Update the root: $z_i^1\gets \exp_{z_i^1}(-\eta G_i)$
\EndIf
\For{$\ell=2,\ldots,m$}
    \State Enforce the slave constraint: $z_i^\ell\gets M_\ell(z_i^1)$
\EndFor
\State \Return updated copies $z_i^1,\ldots,z_i^m$
\end{algorithmic}
\end{algorithm}

\paragraph{Hyperbolic convexity check.}
For the cellular examples, the higher-order faces must remain geodesically convex for mean-value coordinates to be well-defined. This is checked with tangent directions of Poincaré geodesics, not Euclidean chord angles. Let $p=(x_p,y_p)$ and $q=(x_q,y_q)$ be two adjacent vertices. If $x_p y_q-y_p x_q=0$, the geodesic is a Euclidean diameter and the tangent at $p$ toward $q$ is
\[
\tau_{p\to q} = \frac{q-p}{\|q-p\|}.
\]
Otherwise the geodesic is the Euclidean circle orthogonal to $\partial\D$ and passing through $p$ and $q$. Its center $c=(c_x,c_y)$ is found by solving
\[
2c\cdot p=\|p\|^2+1,\qquad 2c\cdot q=\|q\|^2+1,
\]
equivalently, for $\Delta=x_p y_q-y_p x_q$,
\[
c_x = \frac{(\|p\|^2+1)y_q-(\|q\|^2+1)y_p}{2\Delta}, \qquad c_y = \frac{x_p(\|q\|^2+1)-x_q(\|p\|^2+1)}{2\Delta}.
\]
The tangent line is perpendicular to the radius $p-c$, so the two possible unit tangents are
\[
\pm\frac{(-(y_p-c_y),\,x_p-c_x)}{\|p-c\|}.
\]
We choose the sign for which $p+\epsilon\tau_{p\to q}$ is closer to $q$ than $p-\epsilon\tau_{p\to q}$ for a small $\epsilon>0$. For a face $(v_0,\ldots,v_{\ell-1})$, define
\[
a_k=-\tau_{v_k\to v_{k-1}}, \qquad b_k=\tau_{v_k\to v_{k+1}}.
\]
The signed turn at $v_k$ is the oriented angle from $a_k$ to $b_k$:
\[
\theta_k = \Arg\!\left( \left\langle a_k,b_k\right\rangle + \sqrt{-1}\,\det(a_k,b_k) \right)\in(-\pi,\pi].
\]
The face is accepted as geodesically convex when all $\theta_k$ have the sign of the face orientation. The minimized cellular endpoints used below passed this tangent-angle convexity test and had no Poincaré-geodesic edge crossings.

\paragraph{Mean-value directed weights.}
After computing the two endpoint convex embeddings, we discard the scalar endpoint weights and recompute directed mean-value edge weights. Around a center vertex $i$, all neighbors are first attached into the same local star using the side-pairing Möbius transformations.

The star is moved by a disk isometry sending $z_i$ to the origin. For a neighbor occurrence $j$, let
\begin{equation}
\label{eq:local-star-coordinate}
u_j=\frac{z_j-z_i}{1-\overline{z_i}z_j} = r_j e^{\sqrt{-1}\phi_j}, \qquad \rho_j=2\operatorname{arctanh} r_j.
\end{equation}
After sorting the neighbor occurrences by angle, let $\alpha_j^{-}$ and $\alpha_j^{+}$ be the adjacent angular gaps before and after occurrence $j$. The unnormalized hyperbolic mean-value directed weight is
\begin{equation}
\label{eq:directed-mean-value-weight}
\omega_{ij} = \frac{ \tan(\alpha_j^{-}/2)+\tan(\alpha_j^{+}/2) }{ \sinh \rho_j }.
\end{equation}
When multiple attached occurrences have the same angular direction, the numerator is shared equally among them.

\begin{algorithm}[H]
\caption{Directed mean-value weights on the attached local star}
\label{alg:directed-mean-value-weights}
\begin{algorithmic}[1]
\State \textbf{Input:} Graph $G=(V,E)$, position array $Z$, side-pairing maps
\State \textbf{Output:} Directed weights $\omega_{ij}$ for all oriented edge incidences
\ForAll{root vertices $i\in V$}
    \State Build the attached local star of $i$ using $Z$ and the side-pairing maps
    \ForAll{attached neighbor occurrences $j$}
        \State Compute $u_j,\phi_j,\rho_j$ by Eq.~\eqref{eq:local-star-coordinate}
    \EndFor
    \State Sort by $\phi_j$ and compute the adjacent gaps $\alpha_j^{-},\alpha_j^{+}$
    \ForAll{attached neighbor occurrences $j$}
        \State Set $\omega_{ij}$ by Eq.~\eqref{eq:directed-mean-value-weight}
        \State Split the numerator among coincident directions, if any
    \EndFor
\EndFor
\State \Return all directed weights $\omega_{ij}$ and $\omega_{ji}$
\end{algorithmic}
\end{algorithm}

\paragraph{Morphing solve.}
Let $\omega_{ij}^{0}$ denote the directed mean-value weight from vertex $i$ to vertex $j$ in the unit-weight endpoint, and let $\omega_{ij}^{1}$ denote the corresponding directed weight in the non-uniform endpoint. For frame $k=0,\ldots,24$, with $t_k=k/24$, we linearly interpolate the directed weights,
\[
\omega_{ij}(t_k) = (1-t_k)\,\omega_{ij}^{0} + t_k\,\omega_{ij}^{1}.
\]
The scalar edge coefficient recorded for the frame energy is $a_{ij}(t_k)=(\omega_{ij}(t_k)+\omega_{ji}(t_k))/2$. Each frame is obtained by solving the hyperbolic harmonic-map equations with these interpolated directed weights. The solver state is the full position array $Z$, which contains both root and slave entries. During one solver sweep, Algorithm~\ref{alg:root-slave-gradient-accumulation} is applied to every root. For an interior root, the algorithm sees only the root copy. For a boundary root, it also sees the attached slave copies. For a fixed root, the accumulated gradient is computed but ignored by the update step. After each root update, the corresponding slave entries in $Z$ are overwritten by the attaching maps. The endpoint arrays and every stored frame are assumed to satisfy these attaching constraints. The previous frame is used as the initial iterate for the next frame.

\begin{algorithm}[H]
\caption{Mean-value-coordinate morphing}
\label{alg:mean-value-morph}
\begin{algorithmic}[1]
\State \textbf{Input:} Graph $G=(V,E)$, frame count $N$, tolerance $\varepsilon$, and endpoint positions $Z^0$ and $Z^1$
\State \textbf{Output:} Position frames $Z_0,\ldots,Z_{N-1}$
\State Compute directed weights $\omega^0$ from $Z^0$ by Algorithm~\ref{alg:directed-mean-value-weights}
\State Compute directed weights $\omega^1$ from $Z^1$ by Algorithm~\ref{alg:directed-mean-value-weights}
\State $Z_{-1}\gets Z^0$
\For{$k=0,\ldots,N-1$}
    \State $t_k\gets k/(N-1)$
    \ForAll{directed edges $(i,j)$}
        \State $\omega_{ij}(t_k)\gets (1-t_k)\omega_{ij}^0+t_k\omega_{ij}^1$
    \EndFor
    \State Initialize the frame solve with $Z\gets Z_{k-1}$
    \State $\delta\gets\infty$
    \Repeat
        \State $Z_{\mathrm{old}}\gets Z$
        \ForAll{root vertices $i$}
            \State Apply Algorithm~\ref{alg:root-slave-gradient-accumulation} to the root copy and all slave copies of $i$ using $\omega(t_k)$
            \State Store the returned root and slave positions in $Z$
        \EndFor
        \State $\delta\gets\max_i d_{\D}(Z_i,Z_{\mathrm{old},i})$
    \Until{$\delta<\varepsilon$}
    \State Store $Z_k\gets Z$
\EndFor
\State \Return $Z_0,\ldots,Z_{N-1}$
\end{algorithmic}
\end{algorithm}

All four morphs were generated with 25 frames. The figures below show the evenly spaced frame indices $[0,3,6,9,12,15,18,21,24]$. Green curves are graph edges, the dashed black curve is the fundamental-domain boundary, red vertices are fixed polygon corners, and orange and green faces indicate quadrilateral and pentagonal cells, respectively.

\begin{figure}[t]
\centering
\includesvg[width=\linewidth]{genus2_triangulation_morph_grid}
\caption{Genus--2 triangulation morph. The graph has 190 vertices, 527 edges, and 338 triangular faces. The frames show the mean-value-coordinate morph from the unit-weight endpoint to the non-uniform endpoint.}
\label{fig:g2-triangulation-morph}
\end{figure}

\begin{figure}[t]
\centering
\includesvg[width=\linewidth]{genus2_cells_morph_grid}
\caption{Genus--2 cellular morph. The graph has 190 vertices, 490 edges, 270 triangles, 25 quadrilaterals, and 6 pentagons. The displayed frames preserve geodesic convexity for all quadrilateral and pentagonal cells.}
\label{fig:g2-cells-morph}
\end{figure}

\begin{figure}[t]
\centering
\includesvg[width=\linewidth]{genus3_triangulation_morph_grid}
\caption{Genus--3 triangulation morph on the Klein quartic fundamental domain. The graph has 250 vertices, 677 edges, and 428 triangular faces. The frames show the mean-value-coordinate morph from the unit-weight endpoint to the non-uniform endpoint.}
\label{fig:g3-triangulation-morph}
\end{figure}

\begin{figure}[t]
\centering
\includesvg[width=\linewidth]{genus3_cells_morph_grid}
\caption{Genus--3 cellular morph. The graph has 250 vertices, 557 edges, 208 triangles, 80 quadrilaterals, and 20 pentagons. As in the genus--2 cellular example, mean-value weights are computed from the completed attached local stars, and the displayed frames preserve geodesic convexity for all higher-order cells.}
\label{fig:g3-cells-morph}
\end{figure}

Across all four examples, every generated frame remained embedded in the Poincaré disk. No frame contained crossing Poincaré geodesic edges. For the two cellular morphs, every displayed and generated frame also passed the tangent-angle convexity check on all quadrilateral and pentagonal faces. These experiments indicate that the same endpoint minimization and directed mean-value interpolation procedure applies to triangulations and to cellular decompositions with convex higher-order faces, on both the genus--2 Bolza surface and the genus--3 Klein quartic example.

\end{document}
